\section{パップス＝ギュルダンの定理}\label{節:パップスの定理}

回転体の体積を円筒殻による方法で求めると,軸からの距離に帯の面積を掛けた積分が現れた.一方,重心も座標に帯の面積を掛けた積分で求めた.この2つは同じ積分ではないだろうか.そうなら,回転体の体積は面積と重心だけで書けるはずである.

回転軸を $x$ 軸とし,図形が $y\ge0$ の側にあるとする.図形を $x$ 軸に平行な横の細い帯に分け,高さ $y$ の帯の横の長さを $h(y)$ とする.帯が複数に分かれるなら,\,$h(y)$ はそれらの長さの和とする.図形の高さの範囲を $c\le y\le d$ としよう.

横の帯を $x$ 軸のまわりに回すと,半径 $y$,軸に沿う長さ $h(y)$ の円筒の層になる.式\eqref{式:円筒殻の体積}で $x$ と $y$ の役割を入れ替えれば
\[
V=2\pi\int_c^d y\,h(y)\dd y
\]
である.他方,横の帯の中心の高さは $y$ であるから,面積 $A$ と面積重心の高さ $\overline y$ は
\[
A=\int_c^d h(y)\dd y,\qquad
A\overline y=\int_c^d y\,h(y)\dd y
\]
を満たす.

\sectionref{節:面積重心と線重心}では,\,$A\overline y$ を縦の帯から $\frac12\int_a^b(f^2-g^2)\dd x$ として求めた.横の帯から出した上の積分と同じ値になるのだろうか.図形を縦横の線で細かい長方形に分け,各長方形の面積 $\Delta x\Delta y$ に高さ $y$ を掛けた和 $\sum y\,\Delta x\Delta y$ を考えよう.同じ縦の列にある長方形を先に足すと,列1つ分は $\int_{g(x)}^{f(x)}y\dd y\cdot\Delta x=\frac12\{f(x)^2-g(x)^2\}\Delta x$ に近く,これは縦の帯1本分の寄与である.同じ横の行を先に足すと,行1つ分は $y\,h(y)\Delta y$ に近く,これは横の帯1本分の寄与である.足す順序を変えても和は変わらないので,分割を細かくした極限で2つの積分は一致する.どちらも,図形全体で(高さ)$\times$(面積)を足した同じ量を,分け方を変えて計算しているのだ.厳密な証明には極限の扱いに注意が要るが,ここでは立ち入らない.実際,半径 $a$ の上半円板では $h(y)=2\sqrt{a^2-y^2}$ であり,
\[
\int_0^a 2y\sqrt{a^2-y^2}\dd y
 =\left[-\frac23(a^2-y^2)^{\frac32}\right]_0^a
 =\frac{2a^3}{3}
\]
となって,\sectionref{節:面積重心と線重心}で縦の帯から求めた値と一致する.

$V$ の式と $A\overline y$ の式に,同じ積分 $\int_c^d y\,h(y)\dd y$ が現れている.これを $A\overline y$ に置き換えれば,
\[
V=2\pi\overline yA
\]
を得るのだ.

\begin{theorem}{パップス＝ギュルダンの体積に関する定理}
一様な平面図形の面積を $A$,その面積重心と,図形と同じ平面内にある回転軸との距離を $R$ とする.図形が軸の片側にあり,回転によって内部を重複して数えない場合,できる回転体の体積は
\[
V=2\pi R A
\]
である.すなわち,面積に,重心が1回転するときに進む円周の長さを掛ければよい.
\end{theorem}

上の証明では,境界が区分的に滑らかで,断面積による通常の積分ができる図形を考えている.軸が境界に触れても内部を横切らなければ,上の帯による証明がそのまま使える.軸が内部を横切る場合には,軸の両側が同じ場所を回転して重なり得るので,この定理は使えない.重心が軸上にあるから体積が $0$ になる,という使い方はできないのである.

\subsection{回転面の面積}

体積では面積重心が現れた.では,曲線を回してできる面の面積ではどうなるだろうか.結論から言えば,今度は線重心が現れる.その理由を,回転面の小さな部分から考えよう.

回転軸と同じ平面内にあり,軸の片側にある曲線を考える.曲線を短い線分の折れ線で近似し,各線分が回ってできる面の面積を足し合わせて,定積分に帰着させる.

軸からの距離が $r_1,r_2$ の2点を結ぶ長さ $\Delta s$ の線分を軸のまわりに回すと,一般には円すい台の側面ができる.軸に平行な線分なら円筒面,軸に垂直な線分なら円環または円板になる.いずれの場合も,その面積は
\[
\pi(r_1+r_2)\Delta s
\]
である.円すい台の場合は,大きい円すいの側面積から切り取った円すいの側面積を引き,相似比を使って整理するとこの式になる.円筒面では $r_1=r_2$ として $2\pi r_1\Delta s$,円環では $\Delta s=|r_2-r_1|$ として $\pi|r_2^2-r_1^2|$ であり,どちらもこの式に一致する.線分を十分短くすれば $r_1,r_2$ はほぼ同じ値 $r$ になるので,面積はほぼ $2\pi r\Delta s$ である.

曲線が区分的に連続微分可能で,回転によって同じ面を重複して作らないとき,折れ線近似の極限として,回転面の面積は
\begin{equation}\label{式:回転面の面積}
S=2\pi\int_C r\dd s
\end{equation}
になる.ここでは,この極限の正当化の詳細は省く.たとえば $y=f(x)\geq0$($a\leq x\leq b$)を $x$ 軸のまわりに回すなら,\,$r=f(x)$,\,$\dd s=\sqrt{1+f'(x)^2}\dd x$ であるから
\[
S=2\pi\int_a^b f(x)\sqrt{1+f'(x)^2}\dd x
\]
となる.得られるのは,指定した曲線が作る面の面積である.立体の両端を塞ぐ円板などは,その曲線に含めていなければ別に数える.

\begin{theorem}{パップス＝ギュルダンの面積に関する定理}
一様な平面曲線の長さを $L$,その線重心と,曲線と同じ平面内にある回転軸との距離を $R$ とする.曲線が軸の片側にあり,1回転によって同じ面を重複して作らない場合,曲線が作る回転面の面積は
\[
S=2\pi R L
\]
である.
\end{theorem}

実際,軸を $x$ 軸,曲線を $y\geq0$ の側にとれば $r=y$ である.式\eqref{式:線の重心}より $\int_C y\dd s=L\overline y=LR$ となるので,式\eqref{式:回転面の面積}から $S=2\pi RL$ が従う.回転軸の位置や向きが違っていても,軸に合わせて座標をとり直せば同じ議論になる.

2つの定理で掛ける円周は,図形の端や図形の中心が描く円周ではなく,それぞれの重心が描く円周である.体積の定理では面積 $A$ と面積重心を,面積の定理では長さ $L$ と線重心を組にして使う.

\begin{bookpoint}
\textbf{POINT}\par
体積は(面積重心が描く円周)$\times$(面積),回転面の面積は(線重心が描く円周)$\times$(長さ)で求める.
\end{bookpoint}

\begin{example}{半円板の回転}
$a>0$,\,$b>0$ とする.上半円板 $x^2+y^2\leq a^2$,\,$y\geq0$ を,直線 $y=-b$ のまわりに1回転してできる立体の体積と全表面積を求めよ.
\end{example}

\paragraphbf{方針}
回転半径は,重心の高さに $b$ を加えたものである.体積には半円板の面積重心を使う.表面積は境界を円弧と直径に分け,それぞれの線重心を使う.

\paragraphbf{解答}
体積 $\pi^2a^2b+\frac{4\pi a^3}{3}$,全表面積 $2\pi^2ab+4\pi a^2+4\pi ab$

\paragraphbf{解法}
図\ref{図:半円の二つの重心}のように,回転軸は図形の外にある.

\textbf{体積}\quad 半円板の面積は $\frac{\pi a^2}{2}$,面積重心の高さは式\eqref{式:半円板の重心}より $\frac{4a}{3\pi}$ である.回転半径は $b+\frac{4a}{3\pi}$ なので,
\begin{align*}
V&=2\pi\left(b+\frac{4a}{3\pi}\right)\cdot\frac{\pi a^2}{2}\\
 &=\pi^2a^2b+\frac{4\pi a^3}{3}.
\end{align*}
断面でも確かめておこう.位置 $x$ の断面は,外半径 $b+\sqrt{a^2-x^2}$,内半径 $b$ の円環である.よって
\begin{align*}
V&=\pi\int_{-a}^a\left\{\left(b+\sqrt{a^2-x^2}\right)^2-b^2\right\}\dd x\\
 &=\pi\int_{-a}^a(a^2-x^2)\dd x+2\pi b\int_{-a}^a\sqrt{a^2-x^2}\dd x
\end{align*}
となる.第1項は $\pi\cdot\frac{4a^3}{3}=\frac{4\pi a^3}{3}$ である.第2項の積分は半円板の面積 $\frac{\pi a^2}{2}$ に等しいので,第2項は $\pi^2a^2b$ である.よって結果は一致する.

\begin{figure}[htbp]
\centering
\begin{tikzpicture}[scale=1.25,>=Stealth]
\fill[blue!10](-2,0)arc[start angle=180,end angle=0,radius=2]--cycle;
\draw[blue!75!black,thick](2,0)arc[start angle=0,end angle=180,radius=2];
\draw[thick](-2,0)--(2,0);
\draw[->](-2.5,0)--(2.6,0)node[right]{$x$};
\draw[->](0,-1.5)--(0,2.45)node[above]{$y$};
\draw[orange!80!black,thick](-2.5,-1)--(2.5,-1)node[right]{$y=-b$};
\fill[blue!75!black](0,0.8488)circle(0.045);
\fill[red!75!black](0,1.2732)circle(0.045);
\draw[blue!75!black](0,0.8488)--(2.5,0.6)node[right]{面積重心 $\frac{4a}{3\pi}$};
\draw[red!75!black](0,1.2732)--(2.5,1.6)node[right]{線重心 $\frac{2a}{\pi}$};
\draw[<->](-2.25,-1)--(-2.25,0)node[midway,left]{$b$};
\node[below]at(-2,0){$-a$};\node[below]at(2,0){$a$};
\end{tikzpicture}
\caption{同じ半円でも,板と円弧では重心が違う.回転半径は軸から測る.}\label{図:半円の二つの重心}
\end{figure}

\textbf{全表面積}\quad 立体の表面は,境界の円弧と直径がそれぞれ回ってできる.円弧の長さは $\pi a$,線重心の高さは式\eqref{式:半円弧の重心}より $\frac{2a}{\pi}$ であるから,円弧が作る面の面積は
\begin{align*}
S_{\mathrm{円弧}}
 &=2\pi\left(b+\frac{2a}{\pi}\right)\cdot\pi a\\
 &=2\pi^2ab+4\pi a^2.
\end{align*}
これだけでは全表面積にならない.直径 $-a\leq x\leq a$,\,$y=0$ も,回転すると半径 $b$,軸方向の長さ $2a$ の円筒面を作る.その面積は
\[
S_{\mathrm{直径}}=2\pi b\cdot2a=4\pi ab
\]
である.2つの面は端の円周で接するだけで,面積のある部分では重ならない.よって全表面積は
\[
S=2\pi^2ab+4\pi a^2+4\pi ab
\]
となる.\,$x=\pm a$ では外半径と内半径がともに $b$ で立体の厚みが $0$ になるので,端を塞ぐ円板を加える必要はない.
%
